% Part: computability
% Chapter: recursive-functions
% Section: trees

\documentclass[../../../include/open-logic-section]{subfiles}

\begin{document}

\olfileid{cmp}{rec}{tre}
\olsection{वृक्ष}

ज्याप्रमाणे क्रमिका संख्यांनी आणि त्यांचे गुणधर्म व त्यांवरील क्रिया त्यांच्या
संकेतांकांवरील आदिम पुनरावर्ती संबंधांनी व फलनांनी दर्शवता येतात, त्याचप्रमाणे
वृक्ष नैसर्गिक संख्यांनी दर्शवणे कधीकधी उपयोगी ठरते. हे करण्यासाठी आपण
क्रमिका आणि त्यांचे संकेतांक वापरू. वृक्ष म्हणजे एकच गाठ (जिला नामांकन असू
शकते), किंवा अनेक उपवृक्षांशी जोडलेली गाठ (जिला नामांकन असू शकते). त्या
गाठीला वृक्षाचे \emph{मूळ} म्हणतात आणि तिला जोडलेल्या उपवृक्षांना त्याचे
\emph{तात्काळ उपवृक्ष} म्हणतात.

वृक्षाचा पुनरावर्ती संकेतांक $\tuple{k, d_1, \dots, d_k}$ या क्रमिकेने
देऊ; येथे $k$ ही तात्काळ उपवृक्षांची संख्या आणि $d_1$, \dots,~$d_k$ हे
त्यांच्या संकेतांकांचे मूल्य आहेत. गाठींना नामांकने असतील, तर ती तात्काळ
उपवृक्षांनंतर समाविष्ट करता येतात. म्हणून फक्त $l$ नामांकन असलेली एक गाठ
मिळून बनलेल्या वृक्षाचा संकेतांक $\tuple{0,l}$ असेल; आणि $l_1$ नामांकन
असलेले मूळ हे $l_2$, $l_3$ नामांकने असलेल्या दोन एक-गाठीच्या वृक्षांना
जोडलेले असेल, तर त्या वृक्षाचा संकेतांक $\tuple{2,
  \tuple{0, l_2}, \tuple{0, l_3}, l_1}$ असेल.

\begin{prop}
  \ollabel{prop:subtreeseq}
  संकेतांक~$t$ असलेल्या वृक्षाच्या सर्व उपवृक्षांचे संकेतांक घटक म्हणून
  असलेल्या क्रमिकेचा संकेतांक परत करणारे $\fn{SubtreeSeq}(t)$ हे फलन
  आदिम पुनरावर्ती आहे.
\end{prop}

\begin{proof}
  प्रथम, $\fn{ISubtrees}(t) = \fn{subseq}(t, 1, (t)_0)$ हे आदिम पुनरावर्ती
  फलन आहे आणि वृक्ष~$t$ च्या तात्काळ उपवृक्षांचे संकेतांक परत करते, हे लक्षात
  घ्या. आता मूळापासून जास्तीत जास्त $n$ गाठी दूर असलेल्या सर्व उपवृक्षांची,
  संभाव्य पुनरुक्तींसह, क्रमिका काढणारे साहाय्यक फलन
  $\fn{hSubtreeSeq}(t,n)$ परिभाषित करू. मूळापासून जास्तीत जास्त $0$ गाठी
  दूर असलेल्या~$t$ च्या उपवृक्षांची क्रमिका---दुसऱ्या शब्दांत, $t$ च्या
  मुळापासून सुरू होणारी क्रमिका---म्हणजे फक्त~$t$ असलेली क्रमिका. $t$ च्या
  स्तर~$n+1$ पर्यंतच्या सर्व उपवृक्षांची क्रमिका मिळवण्यासाठी स्तर~$n$
  पर्यंतच्या उपवृक्षांची क्रमिका आणि स्तर~$n$ पर्यंतच्या त्या उपवृक्षांच्या
  तात्काळ उपवृक्षांची क्रमिका यांची जोडणी करतो. ही दुसरी क्रमिका मिळवण्यासाठी,
  $f(x)$ हे क्रमिकांचे संकेतांक परत करणारे आदिम पुनरावर्ती फलन असेल, तर
  $g_f(s, k) = f((s)_0) \concat
  \dots \concat f((s)_{k-1})$ हेही आदिम पुनरावर्ती आहे, हे लक्षात घ्या:
    \begin{align*}
      g(s, 0) & = \emptyseq\\
      g(s, k+1) & = g(s, k) \concat f((s)_k)
    \end{align*}
    उदाहरणार्थ, $s$ ही वृक्षांची क्रमिका असेल, तर
    $h(s) = g_{\fn{ISubtrees}}(s, \len{s})$ ही~$s$ च्या घटकांच्या सर्व
    तात्काळ उपवृक्षांची क्रमिका देते. तिच्या साहाय्याने
    $\fn{hSubtreeSeq}$ ची पुढीलप्रमाणे व्याख्या करता येते:
    \begin{align*}
      \fn{hSubtreeSeq}(t, 0) & = \tuple{t} \\
      \fn{hSubtreeSeq}(t, n+1) & = \fn{hSubtreeSeq}(t, n) \concat
      h(\fn{hSubtreeSeq}(t, n)).
    \end{align*}
    संकेतांक~$t$ असलेल्या वृक्षातील उपवृक्षांचा महत्तम स्तर, म्हणजे मूळ आणि
    एखादी पर्णगाठ यांमधील महत्तम अंतर, संकेतांक~$t$ ने परिबद्ध असतो. म्हणून
    संकेतांक~$t$ असलेल्या वृक्षाच्या सर्व उपवृक्षांच्या संकेतांकांची क्रमिका
    $\fn{hSubtreeSeq}(t, t)$ ने मिळते.
\end{proof}

\begin{prob}
  \olref[cmp][rec][tre]{prop:subtreeseq} च्या सिद्धतेतील
  $\fn{hSubtreeSeq}$ च्या व्याख्येत सर्वसाधारणपणे पुनरुक्ती येतात. परिणामी
  क्रमिकेत प्रत्येक उपवृक्षाचा संकेतांक फक्त एकदाच येईल याची खात्री करणारी
  पर्यायी व्याख्या द्या.
\end{prob}

\end{document}
